%% Archived from sections/approximations.tex: finite-state kinetic stability.
%% Removed because it is conditional machinery that cannot be validated by
%% anything else in the paper. Restore by re-appending to approximations.tex.

\subsection{Finite-state kinetic stability with a complete stated generator}
Let $S=\{1,\ldots,m\}$ be a nonempty finite state space. A row-generator
$Q$ satisfies $Q_{ij}\ge0$ for $i\ne j$ and
$Q_{ii}=-\sum_{j\ne i}Q_{ij}$. Row distributions evolve as
$p(t)=p(0)e^{tQ}$. The following result compares two fully specified,
time-homogeneous generators on the same state space. It makes no
irreducibility or reversibility assumption.

\begin{lemma}[Stochastic semigroups and contraction]
\label{lem:markov}
For such a generator, $P_t=e^{tQ}$ has nonnegative entries and every row
sums to one for $t\ge0$. For any real row vector $z$,
$\norm{zP_t}_1\le\norm z_1$.
\end{lemma}
\begin{proof}
Choose a real $\nu>0$ with $\nu\ge\max_i(-Q_{ii})$ and put
$P=I+Q/\nu$. Off-diagonal entries and diagonal entries of $P$ are
nonnegative and its rows sum to one. Products of such matrices have the
same properties. Since $Q=\nu(P-I)$ and $I$ commutes with $P$,
\[
 e^{tQ}=e^{-\nu t}\sum_{n=0}^\infty\frac{(\nu t)^n}{n!}P^n.
\]
The coefficients are nonnegative and sum to one. Entrywise convergence
therefore proves the stochastic property. Finally,
\[
 \sum_j\left|\sum_i z_i(P_t)_{ij}\right|
 \le\sum_i|z_i|\sum_j(P_t)_{ij}=\sum_i|z_i|,
\]
which is the asserted contraction.
\end{proof}

\begin{theorem}[Finite-state generator perturbation bound]
\label{thm:ctmc}
Let $Q,\widehat Q$ be row-generators on $S$ and define
\begin{equation}
 \eta=\frac12\max_i\sum_{j=1}^m|Q_{ij}-\widehat Q_{ij}|.
 \label{eq:eta}
\end{equation}
The sum includes the diagonal entry. For the same initial probability
row vector $p$, put $p_t=pe^{tQ}$ and
$\widehat p_t=pe^{t\widehat Q}$. Then, for every $t\ge0$,
\begin{equation}
 \TV(p_t,\widehat p_t):=\tfrac12\norm{p_t-\widehat p_t}_1
 \le\min(1,t\eta).
 \label{eq:ctmc-tv}
\end{equation}
For any real observable $f:S\to\R$, define
$\operatorname{range}(f)=\max_i f(i)-\min_i f(i)$. Then
\begin{equation}
 |p_t f-\widehat p_t f|
 \le\operatorname{range}(f)\min(1,t\eta).
 \label{eq:observable}
\end{equation}
If the initial distributions differ, the corresponding bound is
$\min(1,\TV(p,\widehat p)+t\eta)$.
\end{theorem}
\begin{proof}
All matrix exponentials are differentiable in finite dimension. For a fixed
$t$, differentiate $F(s)=e^{sQ}e^{(t-s)\widehat Q}$ on $[0,t]$.
Commutation of a matrix with its own exponential gives
\[
 F'(s)=e^{sQ}(Q-\widehat Q)e^{(t-s)\widehat Q}.
\]
Integrating from $0$ to $t$ yields the Duhamel identity
\begin{equation}
 e^{tQ}-e^{t\widehat Q}
 =\int_0^t e^{sQ}(Q-\widehat Q)e^{(t-s)\widehat Q}\,ds.
 \label{eq:duhamel}
\end{equation}
For every probability vector $v$, the triangle inequality gives
\[
 \norm{v(Q-\widehat Q)}_1
 \le\sum_i v_i\sum_j|Q_{ij}-\widehat Q_{ij}|\le2\eta.
\]
By Lemma~\ref{lem:markov}, $pe^{sQ}$ is a probability vector, and
right multiplication by $e^{(t-s)\widehat Q}$ is an $\ell^1$
contraction. Left-multiplying \eqref{eq:duhamel} by $p$ and applying
these two estimates under the integral gives
$\norm{p_t-\widehat p_t}_1\le\int_0^t2\eta\,ds=2t\eta$.
Two probability vectors have $\ell^1$ distance at most two, proving
\eqref{eq:ctmc-tv}.

For the observable inequality, put $z=p_t-\widehat p_t$ and
$c=(\max_i f(i)+\min_i f(i))/2$. Since $\sum_i z_i=0$,
\[
 |zf|=\left|\sum_i z_i(f(i)-c)\right|
 \le\frac{\operatorname{range}(f)}2\sum_i|z_i|
 =\operatorname{range}(f)\TV(p_t,\widehat p_t).
\]
Combining with \eqref{eq:ctmc-tv} proves \eqref{eq:observable}.
For different initial distributions, write
\[
 pe^{tQ}-\widehat p e^{t\widehat Q}
 =p(e^{tQ}-e^{t\widehat Q})
   +(p-\widehat p)e^{t\widehat Q}.
\]
The first term has norm at most $2t\eta$, and contraction bounds the
second by $\norm{p-\widehat p}_1$. Division by two and the universal
probability-distance bound complete the proof.
\end{proof}

\begin{corollary}[From individual rate intervals to a generator bound]
\label{cor:rates-generator}
Suppose all off-diagonal entries are specified and have validated intervals
$Q_{ij}\in[l_{ij},u_{ij}]$ with $0\le l_{ij}\le u_{ij}$.
Choose $\widehat Q_{ij}\in[l_{ij},u_{ij}]$ and set both diagonals by
zero row sums. Define
$\epsilon_{ij}=\max(\widehat Q_{ij}-l_{ij},u_{ij}-\widehat Q_{ij})$
for $i\ne j$. Then
\begin{equation}
 \eta\le\max_i\sum_{j\ne i}\epsilon_{ij}.
 \label{eq:rate-generator}
\end{equation}
\end{corollary}
\begin{proof}
Each off-diagonal difference is at most $\epsilon_{ij}$ in absolute
value, and
$|Q_{ii}-\widehat Q_{ii}|=
|\sum_{j\ne i}(Q_{ij}-\widehat Q_{ij})|
\le\sum_{j\ne i}\epsilon_{ij}$.
The full row sum of absolute differences is therefore at most twice this
sum. Substitution in \eqref{eq:eta} proves the claim.
\end{proof}

An entry may be the sum of several event rates; its interval must include
all those contributions. A nominally absent transition is justified only if
its rate is zero by the stated model or is bounded within the error budget.
The theorem does not certify that an event catalog contains every relevant
physical reaction. It also does not compare generators on different state
spaces without an additional embedding or coarse-graining argument. Pulsed
operating conditions and transport can require time-dependent generators or
a larger model, which must be analyzed separately.

A bounded terminal occupancy is an observable covered by
\eqref{eq:observable}. Growth per cycle is not automatically such an
observable: it may involve accumulated event counts, changing surface states,
transport, and a conversion from deposited material to thickness. A pathwise
count is not generally bounded by the range of a state function on $S$.
Thus the CTMC result is a conditional finite-state stability theorem, not a
direct growth-per-cycle guarantee. Establishing the physical network and the
observable mapping remains a separate task.
